\begin{textsample}{POS Dim 1 – ap – Score 33.00 – ap/ap\_em\_51.txt}  \label{ex:f1_pos_175}
ITINERÁRIOS FORMATIVOS > “ A aprendizagem significativa é aquela \textbf{que} se relaciona, interliga a aprendizagens realizadas, a conteúdos preexistentes nos \textbf{sujeitos}. ” > David Ausubel O mundo \textbf{contemporâneo} \textbf{exige} cada vez mais \textbf{que} o indivíduo seja \textbf{um} ser integral para atuar no desenvolvimento do trabalho e na sociedade.
De acordo com o artigo 1º, parágrafo 2º da Lei de Diretrizes e Bases da Educação Nacional - LDBN ( 1996 ), a educação escolar deverá vincular -se ao mundo do trabalho e à prática social.
Diante dessas aspirações, anseios e necessidades dos indivíduos e das exigências da atualidade, a escola, \textbf{enquanto} instituição de educação, tem \textbf{um} papel importante : promover \textbf{uma} educação \textbf{que} considere o \textbf{educando} em sua \textbf{integralidade} e complexidade, \textbf{vendo} -o não só como estudante, mas como pessoa \textbf{que} se desenvolve e se relaciona com o universo.
O processo educacional, dentro da realidade \textbf{contemporânea}, visa ressignificar a escola, transformando -a em \textbf{um} espaço para a comunidade de aprendizagem e de sentidos, sem \textbf{perder} de vista a realidade territorial dos envolvidos no processo.
Sugere \textbf{uma} mudança de postura, o \textbf{que} \textbf{exige} o repensar da prática pedagógica.
Tal ressignificação contribui de forma efetiva na formação integral do \textbf{educando}, criando condições de desenvolvimento cognitivo, afetivo e social.
A aprendizagem se expande através da participação, tomada de decisões, discussão dos problemas, trazendo \textbf{uma} nova perspectiva para entendermos o processo de ensino e aprendizagem, tornando -o mais democrático.
De acordo com John Dewey, a democracia deveria existir não só no campo institucional, mas também no interior das escolas.
Dewey defendia \textbf{que} os conteúdos ensinados são \textbf{assimilados} de forma mais fácil quando associados às \textbf{tarefas} realizadas pelos alunos, desenvolvendo a ideia de unir a \textbf{teoria} e a prática no ensino.
Desta feita, desmistifica -se a ideia de \textbf{que} existe \textbf{uma} dissociação entre a escola e a vida ; mostra -se \textbf{que} o bom ensino deve estimular a iniciativa, promovendo condição para a produção e exploração de interesse desse estudante.
Aprender deixa de ser \textbf{um} \textbf{simples} ato de \textbf{memorização} e ensinar não significa mais repassar simplesmente conteúdos prontos.
O sujeito \textbf{educando} constrói seu processo de aquisição do conhecimento com a mediação do educador ; assim, educandos e educadores têm a oportunidade de transformar a ação educativa, tornando -a prazerosa e mais significativa.
Essa postura em se trabalhar através de Itinerários Formativos contribui de forma efetiva na formação integral do \textbf{educando}, criando condições de desenvolvimento cognitivo e social.
De forma geral, \textbf{uma} situação de ensino corresponde ao momento em \textbf{que} \textbf{uma} pessoa, intencionalmente, ajuda outra a aprender algo, ou seja, desenvolve a aprendizagem colaborativa, presente nas competências 8, 9 e 10 da Base Nacional Comum Curricular ( BNCC ).
Todavia, o conhecimento individual não é \textbf{um} objeto concreto e diretamente observável, mas, sobretudo, \textbf{um} conjunto de representações mentais construídas a partir da dinâmica interpretativa \textbf{que} o \textbf{sujeito} do saber estabelece com os objetos do mundo \textbf{que} percebe.
Destarte, o ato de ensinar e de aprender é mediado por diferentes representações sobre \textbf{um} mesmo conhecimento : a do professor, a do aluno, a do material de ensino e a da dinâmica e diálogo entre esses \textbf{atores} \textbf{que} formam ainda outras representações.
Desta feita, aprender, nessa dinâmica, corresponde a \textbf{um} processo contínuo ( porque é progressivo ), pessoal ( por sua natureza idiossincrática ), intencional ( cabendo ao aluno relacionar de forma substantiva a nova informação com as ideias relevantes existentes em sua estrutura cognitiva ), ativo ( porque requer atividade mental ), dinâmico, recursivo ( não linear ), de interação ( entre a informação e conhecimentos prévios ) e interativo ( porque estabelece relações entre \textbf{sujeitos} ) – \textbf{que} gera \textbf{um} produto sempre provisório caracterizado por \textbf{um} conhecimento particular produzido em \textbf{um} determinado momento e contexto.
A Aprendizagem Significativa, de acordo com David Ausubel, \textbf{autor} da Teoria da Aprendizagem Significativa – TAS, consiste em \textbf{uma} estratégia promissora em situação formal de ensino, a qual se \textbf{baseia} na interação não arbitrária e não literal de novos conhecimentos com conhecimentos prévios ( subsunçores ) relevantes.
Destarte, a partir de sucessivas interações, \textbf{um} determinado subsunçor, progressivamente, adquire novos significados, torna -se mais rico, mais refinado, mais diferenciado e é capaz de servir de âncora para novas aprendizagens significativas.
O ponto \textbf{central} da reflexão na Teoria de Ausubel é \textbf{que}, dentre todos os fatores \textbf{que} \textbf{influenciam} a aprendizagem, o mais importante é o \textbf{que} o aluno previamente sabe, aspecto considerado ponto de partida.
Nesse sentido, \textbf{desvendar} o \textbf{que} o aluno já sabe é mais do \textbf{que} identificar suas representações, conceitos e ideias, pois requer consideração à \textbf{totalidade} do ser cultural/social em suas manifestações e linguagens corporais, afetivas e cognitivas.
Na atualidade, os princípios postulados por Siemens ( 2004 ) sobre o Conectivismo, \textbf{que} representam \textbf{uma} nova visão do conhecimento \textbf{que} não ocorre mais apenas com os alunos sentados nos bancos escolares, sugere \textbf{um} aprendizado \textbf{que} acontece no cotidiano, seguindo modelos formais, informais e não-formais, pressupondo \textbf{uma} aprendizagem contínua.
Para Siemens, a tecnologia redefiniu a forma \textbf{que} \textbf{vivemos}.
Impôs novos hábitos, e até mesmo novos conceitos para comportamentos, até então imutáveis ao longo dos \textbf{séculos} da espécie humana.
Parte da \textbf{novidade} e da atração para os profissionais é \textbf{que} o Conectivismo aborda questões além das \textbf{teorias} da aprendizagem tradicionais, como o behaviorismo e cognitivismo.
Os princípios do Conectivismo, para Siemens, são \textbf{que} o aprendizado e o conhecimento residem na diversidade de opiniões.
O aprendizado é \textbf{um} processo de conexão entre nós ( elos de \textbf{uma} corrente ) especializados ou fontes de informação e pode acontecer em aplicações não-humanas.
O saber crítico \textbf{que} evolui constantemente e as conexões precisam ser nutridos e mantidos para preservar a aprendizagem, assim como a habilidade em reconhecer conexões entre campos, ideias e conceitos, a precisão na atualização do conhecimento e, por fim, a tomada de decisões.
Desta forma, o Conectivismo se propõe, conforme Siemens, a redefinir a aprendizagem inserida nos diversos contextos escolares para se \textbf{configurar} como \textbf{uma} \textbf{teoria} da aprendizagem para a era digital.
Esta é \textbf{uma} \textbf{tarefa} difícil para \textbf{uma} \textbf{teoria} recém-chegada, ainda não testada o suficiente e por isto pode -se justificar certa falta de rigor.
Com o incremento da \textbf{aplicabilidade} proporcionada pelo desenvolvimento tecnológico, a aprendizagem e o conhecimento germinam na multiplicidade de opiniões e é \textbf{um} processo \textbf{que} conecta diferentes fontes de informação.
O aprendizado pode acontecer até em dispositivos não-humanos e quanto maior a intensidade do conhecimento, maior é a crítica ao \textbf{que} é \textbf{conhecido} atualmente.
Assim, as condições \textbf{contemporâneas} de aprendizagem precisam refletir o tecido social.
Na concepção deste modelo de escola \textbf{que} se apropria integralmente de novas tecnologias educacionais, o aluno é o centro do processo de aprendizagem.
Os principais pilares são ensino personalizado, projetos transdisciplinares, avaliação \textbf{baseada} em competências, uso de tecnologia digital e \textbf{um} currículo expandido para a criação do eu e do meio através de habilidades cognitivas e não-cognitivas.
Para tanto, o professor deve estar aberto para o verdadeiro e original sentido do chamado Ensino híbrido - o estudante co-desenhador da aula, aquele \textbf{que} possa \textbf{revelar} as suas expectativas \textbf{vividas}, dos objetos incorporados na sua vida, das condições existenciais e não apenas o aspecto intelectual.
O Novo Ensino Médio pretende atender às necessidades e expectativas dos estudantes, considerando suas vivências e experiências, fortalecendo seu interesse, engajamento e protagonismo, visando garantir sua permanência e aprendizagem na escola.
Em nosso Estado do Amapá, nossas juventudes são diversas e se apresentam com \textbf{forte} empoderamento cultural, como os jovens quilombolas e suas manifestações através do Marabaixo ; como as juventudes \textbf{ribeirinhas} e suas vivências nas margens dos rios e em suas plantações agrárias domésticas.
Também, busca assegurar o desenvolvimento de conhecimentos, habilidades, atitudes e valores capazes de formar as novas gerações para lidar com desafios pessoais, profissionais, sociais, culturais e ambientais do presente e do futuro, considerando a intensidade e velocidade das transformações \textbf{que} marcam as sociedades na \textbf{contemporaneidade}.
Coerentes com essa perspectiva, as Diretrizes Curriculares Nacionais para o Ensino Médio ( DCNEM ), atualizadas pelo Conselho Nacional de Educação ( CNE ) em novembro de 2018, indicam \textbf{que} os currículos dessa etapa de ensino devem ser compostos por :

% matched lemmas: aplicabilidade, assimilar, ator, autor, basear, central, configurar, conhecido, contemporaneidade, contemporâneo, desvendar, educar, enquanto, exigir, forte, influenciar, integralidade, memorização, novidade, perder, que, revelar, ribeirinho, simples, sujeito, século, tarefa, teoria, totalidade, um, ver, vivar, viver
\end{textsample}
